% Part: first-order-logic
% Chapter: sequent-calculus
% Section: derivations

\documentclass[../../../include/open-logic-section]{subfiles}

\begin{document}

\iftag{FOL}
      {\olfileid[pt]{fol}{seq}{der}}
      {\olfileid[pt]{pl}{seq}{der}}

\olsection{\usetoken{P}{derivation}}

\begin{explain}
Já dissemos como é um sequente inicial e demos as regras de
inferência.  !!^{derivation}s no cálculo de sequentes são
geradas indutivamente a partir disso: cada !!{derivation} ou é um
sequente inicial por si só, ou consiste em uma ou duas !!{derivation}s
seguidas de uma inferência.
\end{explain}

\begin{defn}[$\Log{LK}$-!!{derivation}]
Uma \emph{$\Log{LK}$-!!{derivation}} de um sequente~$S$ é uma árvore finita
de sequentes que satisfaz as seguintes condições:
\begin{enumerate}
\item Os sequentes no topo da árvore são sequentes iniciais.
\item O sequente na base da árvore é~$S$.
\item Todo sequente na árvore exceto $S$ é uma premissa de uma aplicação
  correta de uma regra de inferência cuja conclusão está diretamente
  abaixo daquele sequente na árvore.
\end{enumerate}
Dizemos então que $S$ é o \emph{sequente final} da !!{derivation} e
que $S$ é \emph{!!{derivable} em $\Log{LK}$} (ou $\Log{LK}$-!!{derivable}).
\end{defn}

\begin{ex}
Todo sequente inicial, por exemplo, $!C \Sequent !C$, é !!a{derivation}. Podemos
obter uma nova !!{derivation} a partir disso aplicando, digamos, a regra
$\LeftR{\Weakening}$,
\begin{prooftree}
\Axiom$ \Gamma \fCenter \Delta $
\RightLabel{\LeftR{\Weakening}}
\UnaryInf$ !A, \Gamma \fCenter \Delta$
\end{prooftree}
A regra, contudo, deve ser geral: podemos substituir o $!A$ na
regra por qualquer !!{sentence}, por exemplo, também por~$!D$. Se a premissa
corresponde ao nosso sequente inicial $!C \Sequent !C$, isso significa que tanto
$\Gamma$ quanto $\Delta$ são apenas~$!C$, e a conclusão seria então
$!D, !C \Sequent !C$. Assim, o seguinte é !!a{derivation}:
\begin{prooftree}
\Axiom$ !C \fCenter !C $
\RightLabel{\LeftR{\Weakening}}
\UnaryInf$ !D, !C \fCenter !C$
\end{prooftree}
Podemos agora aplicar outra regra, digamos $\LeftR{\Exchange}$, que nos permite
trocar duas !!{sentence}s à esquerda. Assim, o seguinte também é
uma !!{derivation} correta:
\begin{prooftree}
\Axiom$ !C \fCenter !C $
\RightLabel{\LeftR{\Weakening}}
\UnaryInf$ !D, !C \fCenter !C$
\RightLabel{\LeftR{\Exchange}}
\UnaryInf$ !C, !D \fCenter !C$
\end{prooftree}
Nesta aplicação da regra, que foi dada como
\begin{prooftree}
\Axiom$ \Gamma, !A, !B, \Pi \fCenter \Delta $
\RightLabel{\LeftR{\Exchange}}
\UnaryInf$ \Gamma, !B, !A, \Pi \fCenter \Delta,$
\end{prooftree}
tanto $\Gamma$ quanto $\Pi$ eram vazios, $\Delta$ é $!C$, e os papéis de
$!A$ e $!B$ são desempenhados por $!D$ e~$!C$, respectivamente. De modo muito
semelhante, vemos também que
\begin{prooftree}
\Axiom$ !D \fCenter !D $
\RightLabel{\LeftR{\Weakening}}
\UnaryInf$ !C, !D \fCenter !D$
\end{prooftree}
é !!a{derivation}. Agora podemos pegar essas duas !!{derivation}s e combiná-las
usando $\RightR{\land}$. Aquela regra era
\begin{prooftree}
\Axiom$\Gamma \fCenter \Delta, !A$
\Axiom$ \Gamma \fCenter \Delta, !B$
\RightLabel{\RightR{\land}}
\BinaryInf$ \Gamma \fCenter \Delta, !A \land !B $
\end{prooftree}
No nosso caso, as premissas devem corresponder aos últimos sequentes das
!!{derivation}s que terminam nas premissas. Isso significa que $\Gamma$ é
$!C, !D$, $\Delta$ é vazio, $!A$ é $!C$ e $!B$ é $!D$. Assim, a
conclusão, se a inferência for correta, é $!C, !D \Sequent !C
\land !D$.
\begin{prooftree}
\Axiom$ !C \fCenter !C $
\RightLabel{\LeftR{\Weakening}}
\UnaryInf$ !D, !C \fCenter !C$
\RightLabel{\LeftR{\Exchange}}
\UnaryInf$ !C, !D \fCenter !C$
\Axiom$ !D \fCenter !D $
\RightLabel{\LeftR{\Weakening}}
\UnaryInf$ !C, !D \fCenter !D$
\RightLabel{\RightR{\land}}
\BinaryInf$ !C, !D \fCenter !C \land !D $
\end{prooftree}
É claro que também podemos inverter as premissas; então $!A$
seria $!D$ e $!B$ seria~$!C$. 
\begin{prooftree}
\Axiom$ !D \fCenter !D $
\RightLabel{\LeftR{\Weakening}}
\UnaryInf$ !C, !D \fCenter !D$
\Axiom$ !C \fCenter !C $
\RightLabel{\LeftR{\Weakening}}
\UnaryInf$ !D, !C \fCenter !C$
\RightLabel{\LeftR{\Exchange}}
\UnaryInf$ !C, !D \fCenter !C$
\RightLabel{\RightR{\land}}
\BinaryInf$ !C, !D \fCenter !D \land !C $
\end{prooftree}
\end{ex}

\end{document}
